\documentclass[10pt]{article}
\usepackage[margin=1.7cm]{geometry}
\usepackage{amsmath,amssymb}
\usepackage{booktabs}
\usepackage{graphicx}
\usepackage[T1]{fontenc}
\usepackage{microtype}
\usepackage[colorlinks,linkcolor=black,citecolor=black,urlcolor=blue]{hyperref}
\newcommand{\op}[1]{\texttt{#1}}
\title{Spheroid + ECM --- archived sections}
\author{}\date{}
\setlength{\parindent}{0pt}
\setlength{\parskip}{0.6em}
\begin{document}\maketitle

\noindent\emph{Split out of \texttt{spheroid\_ecm.tex}. These sections describe the model as it stood
before the external audit of \S\,``External audit'' below, and before junction myosin was re-keyed from
length to tension and the membrane made overdamped. They are kept because the measurements in them are
real and because the audit is the work list; they are not kept because the conclusions stand.}

\section*{6.\; Status and the next operator}

$63$ runs in \texttt{log/okuda\_ECM}; \texttt{prototype/ecm/RUNS.md} records which are mutually
comparable, since three quantities were redefined mid-campaign (notably $|J-1|\to\sigma_{\mathrm{vM}}$,
$\sim\!100\times$ in ``strained fraction''). The coupling of \S4 is one-way; \op{ecm\_load\_3d} is the
staggered two-way scheme, implemented and uncalibrated. \op{cell\_polarity} is absent---the epithelium
has no simulated basal side, its cells being wedges from the centre and the monolayer a rendering
convention---as are \op{cell\_extrude\_3d} and MMP degradation.

The clearest gap is in Eq.~\eqref{eq:bond}. With $N$ fixed on a shell whose radius grows $3.6\times$, the
areal density $N/4\pi R^2$ falls $13\times$ and every rest length must creep to $3.6\times$ its original
span: what $\tau_r$ delivers is therefore \emph{dilution with amnesia}, not remodelling, and it is why the
surviving sheet sits permanently near $10\%$ strain. A basement membrane does not thin thirteenfold as an
acinus grows---the epithelium secretes laminin and collagen~IV continuously and thickness is held roughly
constant. The missing operator is \op{basement\_membrane\_secrete}, kind \op{structural}: insert
particles where areal density has fallen below target and bond them in, the sheet's analogue of
\op{divide\_3d}. Plexus supports it directly through \texttt{Level.free\_slots}/\texttt{spawn} with an
oversized reservoir. Only then does ``the membrane survived'' mean what it means biologically---that the
tissue did not outgrow the membrane it was building---rather than that its crosslinks were stretched and
told to forget.

\section{Two things the implementation gets wrong, stated plainly}

\paragraph{The anchor sets where the sheet sits, and only a narrow band of it works.} With critical
damping the sheet stopped tearing, but it then sank \emph{through} the epithelium: the gap to the
recorded surface ran $+0.0040\to-0.0101\to-0.0255$, every particle inside by mid-run. Tracking lag does
not explain it ($2v_a/\sqrt{k}=0.0019$, sixteen times too small). Nor---and this corrects the obvious
guess---does hoop tension: quartering $k_{\text{bond}}$ from $2\times10^5$ to $5\times10^4$ at fixed
$k_{\text{adh}}$ moves the gap by under 5\% ($+0.00185/{+}0.00183/{+}0.00176$) and leaves the strain
alone. The reason is in Eq.~\eqref{eq:bond}'s companion: remodelling fixes the steady-state strain at
$\dot\epsilon\,\tau=0.0027\times60=0.16$ (measured $0.151$--$0.160$ in every run, at every stiffness),
so rest lengths adapt and bond force never becomes large. What sets the gap is $k_{\text{adh}}$ against a
residual inward force; the untested candidate is the interstitial matrix pressing through the shared
grid, which would be the one real matrix-on-membrane effect in the model.

The sweep is narrow. At 200 frames $k_{\text{adh}}=4\times10^4/6\times10^4/8\times10^4$ gives gaps
$-0.0037/-0.0013/+0.0019$; $1\times10^5$ \emph{reverses} to $-0.069$ and $2\times10^5$ tears 94\% of
the crosslinks. A stiffer spring cannot hold less well, so the reversal is a stability threshold rather
than mechanics, and $8\times10^4$ is the usable value---at which 46\% of particles are still inside.
A membrane pinned at fixed angular positions on a surface that triples in radius is being asked to cover
nine times the area with the same particles; the operator that would resolve this, and which the model
does not have, is \op{basement\_membrane\_secrete}.

\paragraph{\op{activity} is degenerate with $\Lambda$.} Actomyosin entered Eq.~\eqref{eq:avm} as the
single scalar $\Lambda$ and now also as a per-junction variable, and the two do not double-count: myosin
\emph{multiplies} the line term, $\Lambda\sum_e m_e\ell_e$, and the recruitment setpoint of
Eq.~\eqref{eq:myo} is written against the \emph{running mean} edge length
$\langle\ell\rangle$, so
\[ m^{ss}_e=a\Bigl(1+\beta\bigl(\ell_e/\langle\ell\rangle-1\bigr)\Bigr) \quad\Longrightarrow\quad \langle m^{ss}\rangle=a . \]
At $a=1$ the mean is pinned to one, $\Lambda$ keeps the meaning it was calibrated with, and myosin is a
pure zero-mean redistribution---which is why the ablation is bit-identical rather than merely close.
But it follows that $a$, which scales every junction uniformly, is \emph{not identifiable against}
$\Lambda$: the monotone radius response over $a\in[0,3]$ ($7.216\to6.138$) is reproduced by sweeping
$\Lambda$ and is a consistency check, not a result. The new parameters are $\beta$, $\tau$ and
$m_{\text{new}}$.

\paragraph{The epithelial surface is not an entity.} It is $R(\hat{\mathbf u},t)$, a $32\times64$ angular
table, read by four operators (\op{cell\_to\_ecm}, \op{cell\_exclude\_3d}, \op{integrin\_adhesion},
\op{basement\_membrane\_seed}). Being a table and not a Level it carries no state and can receive no
delta, which is the structural reason the adhesion of Eq.~\eqref{eq:integrin} is one-way. The Plexus2
form is a Level \op{surface}, one member per outer face with $(\mathbf x,\mathbf v,\hat{\mathbf n},A)$,
written each frame by a \op{broadcast} from the replayed mesh; the four operators then become ordinary
\op{lateral} ones between two Levels, the angular lookup disappears, and---because a Level can be
integrated---the integrin spring becomes mutual. A stronger version is available once the sheet is
anchored, since the membrane particles are themselves a discretisation of that surface and could carry
the contact directly; but runs without a membrane still need a contact path, so \op{surface} is the
general answer and membrane-as-surface the specialisation.

\section{The basement membrane, drawn out}

\begin{figure}[t]\centering
\includegraphics[width=\textwidth]{fig_bm_schematic.pdf}
\caption{\textbf{What the membrane is made of.} Four objects, and the arrows between them are the whole
model. The \emph{epithelium} (cream wedges) is a recording: it advances on a prescribed schedule and
nothing can push it back. Its outer boundary is the \op{surface} Level (orange), one element per
membrane node, each holding a fixed direction $\hat{\mathbf u}_i$ and the radius $R(\hat{\mathbf u}_i,t)$
the tissue reached in that direction. The \emph{basement membrane} is grey \emph{nodes} joined by green
\emph{crosslinks}, sitting a distance $\delta$ outside. Each node is tied to its \emph{own} surface
element by one red \emph{integrin} spring---one per node, to a fixed angular position, which is what
makes it adhesion rather than contact. Outside is the interstitial matrix (purple), which the membrane
never bonds to: they meet only through the shared grid, or in the spring-graph version, not at all.}
\end{figure}

A basement membrane is a thin sheet of protein that an epithelium lays down underneath itself and sits
on. Two networks do the work: \textbf{laminin}, which cells grip, and \textbf{collagen~IV}, which is
crosslinked and carries the load; \textbf{nidogen} staples the two together and \textbf{perlecan} fills
the mesh. We model the load-bearing one. So a node is not a molecule---it is a patch of sheet about a
micron across---and a crosslink is not one collagen~IV protomer but the stiffness of all of them across
that patch.

Three rules, and everything in Fig.~2 follows from them.

\paragraph{1. The crosslinks make it a sheet.} Every pair of neighbouring nodes is a spring with its own
rest length $\ell^0_{ij}$, and its \emph{elongation} is what the colours in Fig.~2 show:
\[ \epsilon_{ij}=\frac{L_{ij}-\ell^0_{ij}}{\ell^0_{ij}},\qquad
   \mathbf f_{ij}=k\,(L_{ij}-\ell^0_{ij})\,\hat{\mathbf d}_{ij},\qquad
   \text{bond dies if } \epsilon_{ij}>\epsilon_{\rm break}. \]
Without them the nodes are dust: independent points that cannot tear, cannot fragment, and cannot be
measured for either.

\paragraph{2. Integrin pins each node to a \emph{direction}, not to the nearest surface.} The anchor is
\[ \mathbf a_i=\mathbf c+\hat{\mathbf u}_i\bigl(R(\hat{\mathbf u}_i,t)+\delta\bigr), \]
with $\hat{\mathbf u}_i$ frozen when the node is laid down. Anchoring to the \emph{nearest} point instead
would let a node slide freely along the surface---that is contact, not adhesion---and the sheet would
never stretch at all. Run~73 is that experiment run on purpose, and its mean elongation is the lowest in
the series outside the slack cases.

\paragraph{3. Remodelling lets the sheet forget.} A crosslink's rest length creeps toward its current
length on a timescale $\tau$, $\dot\ell^0=(L-\ell^0)/\tau$, which gives the relation that governs this
whole series: a sheet stretched at rate $\dot\varepsilon$ settles at
\[ \boxed{\;\langle\epsilon\rangle \simeq \dot\varepsilon\,\tau\;} \]
independent of stiffness and of supply. It is why $k$ can be swept over three orders and change nothing,
and why $\tau$ is the parameter that decides whether the membrane ever approaches failure.

\paragraph{4. Secretion adds material.} A surface growing from $R_0$ to $R_1$ needs $(R_1/R_0)^2$ times
as many nodes to stay the same thickness. Deposition is a fixed fraction $\rho$ of existing material per
frame, so there is a threshold, predictable before any run:
\[ \rho^\ast=\frac{2\ln(R_1/R_0)}{N_{\rm frames}}=\frac{\ln 12.77}{400}=0.0064\ \text{per frame}, \]
measured at $0.006$.

\begin{figure}[t]\centering
\includegraphics[width=\textwidth]{fig_bm_series.pdf}
\caption{\textbf{Seven configurations, same tissue, end state.} Only the membrane is drawn: nodes grey,
crosslinks coloured by elongation with full scale at the break threshold. Reading them against
Eq.~(boxed): \textbf{71} ($\tau=10$) forgets stretch as fast as it arrives and is uniformly slack---the
whole sheet is dark blue at $\langle\epsilon\rangle=0.028$. \textbf{69} ($\tau=60$) is the same network
(180{,}304 against 187{,}781 crosslinks, both fully covering) under four times the tension, $0.104$, and
reads teal-green. \textbf{70} ($\tau=0$, remodelling removed) cannot forget at all: tension accumulates,
p99 reaches $0.263$ against a threshold of $0.35$, and $39\%$ of the crosslinks have already gone
(109{,}798). \textbf{74} lowers the threshold to $0.08$ instead---the mechanics are untouched, but now
the sheet is above it, and $74\%$ of the crosslinks break (47{,}046); the survivors are the slack ones,
which is why its mean elongation reads \emph{lower} than the reference at $0.033$. \textbf{72} is
starved below $\rho^\ast$: 6{,}624 nodes instead of 37{,}427, a sheet stretched thin rather than torn.
\textbf{73} removes the anchor and the sheet slides instead of stretching---and is the only one that
loses coverage, $0.895$, because nothing holds it in place. \textbf{75} is the same membrane on an ovoid
($1.33$), where the equator carries $33\%$ more elongation than the poles.}
\end{figure}

\subsection*{How it is built in Plexus2}

The hierarchy is four sets and one field. \op{cells} and \op{junction} are the epithelium; \op{surface}
is its boundary; \op{basement\_membrane} is the sheet; \op{interstitial\_ecm} is the matrix; and
\op{mpm\_grid} is the field every particle body scatters into, which is how bodies that never touch each
other exchange momentum.

The membrane's operators divide by what they do to the state, which is the distinction Plexus2 draws:
\op{seed\_basement\_membrane} and \op{basement\_membrane\_secrete} are \textbf{structural} (they change
how much material exists); \op{basement\_membrane\_bond} and \op{integrin\_adhesion} are
\textbf{dynamics} (they return a force the engine integrates); \op{basement\_membrane\_remodel} is
\textbf{field}-like (it rewrites rest lengths in place); and \op{basement\_membrane\_bond\_break} is
\textbf{rewire}, because fragmentation \emph{is} a change of the edge set and is registered as one.

Two implementation choices are worth naming because neither is forced by the biology. First, per-bond
state is keyed by the \emph{unordered pair of node indices} rather than stored by position in an array,
so division, T1 and death permute the mesh freely and no topology operator needs to know that crosslinks
exist. Second, the membrane no longer uses MPM: nothing about the sheet came from the continuum solver,
and dropping it costs exactly one thing---momentum exchange with the matrix, measured as an ECM stress
falling from $7.5$ to $0.006$. That coupling has to come back as an explicit operator
(\op{basement\_to\_ecm}), which resolves at the sheet's own length scale instead of the grid's.


\section{The junction network, drawn out}

\begin{figure}[t]\centering
\includegraphics[width=\textwidth]{fig_junction_schematic.pdf}
\caption{\textbf{What a junction is, and the one loop that makes it interesting.} Left: two cells share
an edge, and that edge \emph{is} the junction---it belongs to neither cell, which is why its state is
keyed by the unordered pair of vertices $(i,j)$ rather than stored on a cell. Myosin $m_e$ sits on the
shared edge and sets its tension. Right: the feedback. A junction longer than average recruits more
myosin ($\beta$), which pulls harder, which shortens it. That loop is the only thing here that Okuda's
model cannot already do.}
\end{figure}

\subsection*{The rules}

In the original vertex model, contractility is one number for the whole tissue: every junction pulls
with the same line tension $\Lambda$. We keep $\Lambda$ and give each junction a dimensionless
multiplier $m_e$ on top of it, so the line term of Eq.~\eqref{eq:avm} becomes
\[ \Lambda\sum_e \ell_e \;\longrightarrow\; \Lambda\sum_e m_e\,\ell_e . \]

\paragraph{1. Myosin chases a setpoint that depends on length.} A junction longer than the current
average recruits more; shorter recruits less. Written against the \emph{running} mean $\langle\ell\rangle$,
\[ m^{ss}_e=a\Bigl(1+\beta\bigl(\ell_e/\langle\ell\rangle-1\bigr)\Bigr),\qquad
   \dot m_e=\bigl(m^{ss}_e-m_e\bigr)/\tau_m . \]
Because the reference is the running mean, $\langle m^{ss}\rangle=a$ exactly: the average is pinned and
myosin is a pure \emph{redistribution}. At $a=1$ nothing is double-counted with $\Lambda$, which is why
the ablation is bit-identical rather than merely close.

\paragraph{2. And therefore $a$ is not a new parameter, but $\beta$ is.} Since $a$ multiplies every
junction equally, it does exactly what scaling $\Lambda$ does; the two cannot be told apart from any
observation. $\beta$ can: it changes the \emph{disorder} of junction lengths at constant tissue size.
Sweeping $a$ over $[0.6,1.4]$ traces a line in (radius, junction-length CV) of slope $+0.103$; $\beta$
moves perpendicular to it, $\Delta R=-0.02$ (0.2\%, nothing) against $\Delta\mathrm{CV}=-0.099$, where
that line predicts $-0.002$. Measured:
\begin{center}\small
\begin{tabular}{lccc}\toprule
 & $a$ & $R$ & CV of $\ell_e$ \\ \midrule
reference        & 1.0 & 10.94 & 0.321 \\
$\beta=2$        & 1.0 & 10.95 & \textbf{0.249} \\
$\beta=4$        & 1.0 & 10.92 & \textbf{0.222} \\
blebbistatin-ish & 0.6 & 11.12 & 0.350 \\
hypercontractile & 1.4 & 10.73 & 0.310 \\ \bottomrule
\end{tabular}
\end{center}
$\beta=4$ reaches a CV no value of $a$ in range approaches, at a radius that has not moved.

\paragraph{3. A junction born from a division inherits.} Not every edge that appears is new. A division
inserts a vertex into each of two edges and joins them, so it makes two \emph{kinds}: a split half of a
cut junction, which has a myosin history, and the daughter--daughter interface, which does not. Measured
over 251 frames: of 13{,}750 edges that had no record, 9{,}920 (72\%) were split halves and 9{,}760 of
those (98.4\%) recovered their parent. Only the remaining 28\% take $m_{\rm new}$.

\paragraph{4. $\tau_m$ was chosen, not measured.} Stated plainly because it is the weakest number in
this section: $\tau_m=20$ frames is a default. It is swept nowhere except a determinism check. What it
sets is how far myosin lags the length changes driving it, and the analogous relation for the membrane
($\langle\epsilon\rangle\simeq\dot\varepsilon\tau$) was measurable, so this one should be too.

\subsection*{How it is built in Plexus2}

One operator, \op{junction\_myosin}, of kind \textbf{structural}: it does not return a force, it rewrites
a per-edge quantity that the mechanics then reads in the same frame. It is scheduled immediately before
\op{shape\_energy\_3d} for that reason---placed after, every relaxation would use the previous frame's
myosin.

The design claim worth stating is what is \emph{not} there. Junction is not a Level. Per-edge state is
keyed by the unordered vertex pair, so \op{divide\_3d}, \op{reconnect\_t1\_3d} and any future death
operator permute, lengthen and shorten the half-edge arrays freely: an edge that still exists finds its
own value, one that is new finds its parent or a default, one that is gone is simply never asked. None
of the three topology operators knows that myosin exists, and adding the next per-edge state will not
require editing them either. Promoting junction to a Level would buy per-edge provisioning and cost a
reconciliation operator for every topology change.


\section{External review, August 2026 --- and what to do about it}

An external reviewer with a cell-biology and biophysics background audited the code and this note. The
findings below are theirs; the three marked \textbf{[verified]} I re-checked independently against the
code and the run archive before recording them. Citations are as supplied and are \emph{not} all
independently confirmed; those I could not check are marked \textbf{[unchecked]}.

\subsection*{Verdict}
The engineering is sound and self-critical; the biology is thinner than the prose implies, and three
headline results are properties of the code rather than of the model.

\subsection*{The three most serious findings}

\paragraph{1. The growth gate cannot see stress magnitude. [verified]} The gate compares
$P/p_{99}$ to $p_{\text{half}}=\mathrm{median}(P_{\text{late}})/p_{99}$, so it evaluates
$f(P/\mathrm{median}\,P)$: \emph{multiplying the entire pressure field by any constant leaves it
unchanged}. Every ``the amount does not matter, only the pattern'' control in this campaign is therefore
an identity of the normalisation, not a finding. Compounding this, \op{smooth\_phi\_deg}$=360$ forces the
map axisymmetric before any cell divides, so an oblate spheroid is the only shape available. The
reviewer notes that Hill $n=6$ also sits beyond the one calibrated fit in the literature: Van Liedekerke
et al.\ 2019 fit this exact form to Helmlinger/Montel data, obtained $n\in[1,2]$, and explicitly reject
$n\ge4$. \textbf{[unchecked]} They further note our own primary citation, Helmlinger 1997, reports
suppression accompanied by \emph{decreased apoptosis with no significant change in proliferation}.

\paragraph{2. The membrane series is coupled to nothing, and its adjudicating metric is \texttt{NaN}.
[verified]} Runs 69--74 return \emph{byte-identical} matrix metrics ($\texttt{strained\_frac\_end}
=0.9635714285714285$, identical \texttt{front\_r95} arrays). In \op{membrane\_impl="graph"} the sheet's
MPM operators are stripped, so it exchanges no momentum with the matrix: the seven configurations are a
sheet driven by a prescribed boundary, affecting nothing. Separately, \op{lcc\_end} is \texttt{NaN} in
every race run, because the component fraction is computed only every 40 frames and 403 is not a
multiple of 40 --- so the quantity \S5 claims is reported ``instead of the bond count'' has never been
reported. Worst, run~74's mean bond degree is $z=2.51$ \textbf{[verified]}, below both rigidity
($z\!\approx\!4$) and connectivity percolation, while runs 69--73, 75 sit at $5.4$--$10.0$: that sheet
most likely has no mechanical integrity, and Fig.~2's reading of it as ``the survivors are the slack
ones'' is probably wrong.

\paragraph{3. Both new levels use the wrong state variable or the wrong integrator.} Junction myosin is
driven by \emph{length}, not tension, and the docstring's identification of the two is false in this
energy. The literature has myosin recruited by tension or strain rate, giving a \emph{positive} feedback
that drives T1; ours is a negative feedback on length that homogenises junctions. Worse, in the
biologically correct fast-myosin limit the operator is analytically an edge spring: substituting
$m_e=a(1+\beta(\ell_e/L-1))$ gives $\Lambda a\sum_e[(1-\beta)\ell_e+\beta\ell_e^2/L]$, a line tension
plus a harmonic spring of stiffness $2\Lambda a\beta/L$. Reducing the CV of edge lengths is exactly what
that does, so our one non-degenerate myosin result is guaranteed by construction. The reviewer also
shows the \op{clamp\_min(0)} breaks the $\langle m\rangle=a$ argument at $\beta=4$ (estimated $+16\%$),
which is the row the table rests on. Separately, the membrane is integrated inertially with unit mass,
so the whole critical-damping / tracking-lag / sinking narrative of \S7 is the phenomenology of an
inertia that should not exist at $\mathrm{Re}\sim10^{-10}$.

\subsection*{Other findings worth keeping}
$\bullet$ $\rho^\ast$ is an identity of \op{basement\_membrane\_secrete}, not a prediction --- and it was
computed from the \emph{ungated} tissue ($R_1/R_0=3.57$) while the sweep ran on the gated one, for which
$\rho^\ast=0.0055$, not $0.0064$. $\bullet$ ``Thinning, not rupture'' is imposed: there is no thickness
variable, and remodelling puts rupture out of reach by construction. The literature has BMs opening by
\emph{focal, protease-mediated perforation}, and being maintained or stiffening during growth.
\textbf{[unchecked]} $\bullet$ Our integrin anchors to the collagen~IV network; integrins bind
\emph{laminin}, and \emph{Drosophila} has no collagen-binding integrin at all. $\bullet$ Break strain
$0.35$ is below every tissue-level BM failure measurement cited. $\bullet$ $\tau_m=20$ frames is
$\sim\!3.5$\,h under the only sensible time anchor, against a junctional myosin FRAP $t_{1/2}$ of
$\sim\!6$\,s. $\bullet$ Our citation ``Ku \& Bilder (2023) \emph{Dev.\ Cell} 58:522'' is wrong and must
be corrected or removed. $\bullet$ Our oblateness estimator returns $1.022$ on a perfect sphere, so
several reported effects sit at its own null.

\subsection*{The work list, in order}
\begin{enumerate}\itemsep2pt
\item \textbf{Fix \op{lcc} and re-report 69--77 by connectivity.} Compute the component fraction on the
final frame unconditionally, and record mean degree $z$ alongside every bond count. No new simulation
--- the trajectories are on disk. This settles run~74 and decides ``thinning vs fragmented''.
\item \textbf{Three tissue passes that decide the ovoid:} \op{gate\_hill}$=1$ and $=2$ at fixed floor,
and one run at \op{smooth\_phi\_deg}$=90$ so axisymmetry is not imposed. Then quote the estimator's
sphere null and rewrite the ``amount does not matter'' controls as the normalisation identity they are.
\item \textbf{Give junction myosin the per-edge tension} (available from \op{shape\_energy\_3d}'s
gradient) or the strain rate, report $\langle m\rangle$ against $a$, and measure something
non-tautological --- T1 rate or hexagonal fraction against a $\Lambda$-matched control.
\item \textbf{Make the epithelium feel the sheet}, for runs 70 and 74 only: they are the two
configurations carrying real tension, and a local hernia through a fragmented patch is the one
phenomenon a basement-membrane model exists to produce.
\item \textbf{Switch the membrane to overdamped integration}, which deletes \S7 rather than fixing it.
\end{enumerate}


\section{External audit, and the work list it produces}

An external reviewer with a cell-biology and biophysics background was asked to audit this prototype
adversarially: to find where the biology is wrong, where a simplification changes a conclusion, and where
a reported result is a property of the code rather than of the model. What follows is that audit,
condensed. I have marked \textbf{[verified]} against every claim I checked independently before writing
it down; those checks are cheap and are listed so the reader can repeat them.

\subsection*{Findings I confirmed}

\paragraph{The growth gate cannot see stress magnitude. [verified]} The gate reads
$\mathrm{press}/p_{1/2}$ where the map is normalised as $P/P_{99}$ and $p_{1/2}$ is the late-run median in
those same units. The normalisation cancels: the operator is a function of $P/\mathrm{median}(P)$ alone,
so multiplying the entire pressure field by any constant---stiffer matrix, denser matrix, larger
$k_{\text{contact}}$---leaves the gate unchanged. The comment in \texttt{load\_ops.py} states that the
normalisation was introduced so that $p_{1/2}$ would mean the same thing whatever the matrix stiffness.
That is an invariance built on purpose; \S4's ``five amplifications each raised the stress and lowered the
shape'' and ``the shape comes from the pattern, not the amount'' then report the null of a sweep the code
guarantees to be null. Compounding it, \texttt{smooth\_phi\_deg}$=360$ makes the map axisymmetric before
any cell divides, so the output is constrained to be a spheroid of revolution, and $p_{1/2}$ is
calibrated on the last quarter of the same run it gates from frame~0.

\paragraph{The Hill exponent is outside the range the one calibrated fit allows.} Van Liedekerke
et~al.\ (2019, \emph{PLoS Comput Biol} 15:e1006273) fit this exact functional form to the
Helmlinger/Montel/Alessandri data and obtain $n\in[1,2]$, stating that $n\ge4$ produces a flattening of
residual growth that is not observed. We use $n=6$, and the only sweep that exists (4/6/8) spans the flat
top of the switch and co-varies \texttt{floor} with \texttt{hill}, confounding the two. Integrating the
gate on our own recorded map, the pole/equator contrast falls only about 16\% at $n=1$, so the ovoid may
well survive---but $n=1$ and $n=2$ have never been run on the headline configuration.

\paragraph{Helmlinger measured the opposite of what we cite it for.} Our primary citation for
pressure-suppressed proliferation states in its abstract that stress-induced inhibition of plateau-phase
spheroids ``was accompanied by decreased apoptosis with \emph{no significant changes in
proliferation}''. Delarue et~al.\ (2014) locate the causal variable in cell \emph{volume} acting through
p27$^{\text{Kip1}}$, two abstractions from a Hill function of stress; and the sign is not universal
(Aragona 2013, Gudipaty 2017 have proliferation permitted where stretch is high).

\paragraph{The oblateness estimator's own null is 1.020. [verified]} On 200{,}000 uniformly sampled
points of a perfect sphere, $r_{\text{eq}}/r_{\text{ax}}$ taken as ratio of 98th percentiles reads
\textbf{1.020}, while the RMS form reads 1.000. So the reference tissue's ``aspect 1.021'' is exactly
round, and every effect in the campaign below $\approx\!1.05$---including the myosin runs at
1.016--1.056---sits at or under one estimator's null. Two estimators are in use and neither null was
quoted. This is a fourth quantity redefined mid-campaign and belongs in \texttt{DEFECTS.md}.

\paragraph{The membrane series is a sheet coupled to nothing. [verified]} Runs 69--74 return
\emph{byte-identical} interstitial-matrix metrics---\texttt{strained\_frac\_end} $=0.9635714285714285$,
\texttt{max\_disp} $=0.20738065242767334$, \texttt{front\_reaches\_wall} $=235$---to full float precision
across all six. In \texttt{membrane\_impl="graph"} the membrane's MPM operators are stripped, so the
sheet exchanges no momentum with the matrix at any resolution. The \S5 caveat about $dx$ failing to
resolve a $0.002$ sheet is moot for these runs: there is no coupling to under-resolve.

\paragraph{The fragmentation metric has never returned a number. [verified]} \texttt{lcc\_end} reads
\texttt{bt[-1,3]}, but the component fraction is only computed when the frame index is a multiple of
\texttt{components\_every}$=40$ and is \texttt{NaN} otherwise. With 403 frames it is \texttt{NaN} in all
27 race runs and is not recorded at all for 69--77. \S5 asserts that what is reported is the largest
connected component rather than the bond count. It has never been reported.

\paragraph{Run 74 is probably not a sheet. [verified]} 47{,}046 bonds on 37{,}424 live nodes is a mean
degree of \textbf{2.51}, against $z\approx4$ for central-force rigidity percolation in 2D and
$z\approx4.5$ for connectivity percolation. Figure~2 reads that panel as ``the survivors are the slack
ones''; the more likely reading is that the network has lost mechanical integrity altogether.
\texttt{coverage}$=1.0$ does not contradict this---it counts occupied angular bins, not connectivity.

\paragraph{$\langle\epsilon\rangle\simeq\dot\varepsilon\tau$ fails at the $\tau$ we use.} With
$\dot\varepsilon=0.00314$/frame: $\tau=10$ predicts 0.031 against 0.028 measured (good); $\tau=60$
predicts 0.188 against \textbf{0.104} measured, 44\% low. The values quoted as verification (0.151--0.160)
come from the \emph{non-secreting} runs; in the secreting series most of the final material was deposited
unstrained, which dilutes the mean. The boxed law is presented as holding across the series and does not.

\paragraph{The secretion threshold is an identity, not a prediction.} $\rho^\ast$ follows in three lines
from the operator's own setpoint $N_{\text{want}}=n_0(R/R_0)^2$. Deriving it and then measuring it is a
unit test presented as a falsifiable prediction---and the number quoted uses the \emph{ungated} tissue
($R_1/R_0=3.57$) while the sweep runs on the gated one ($\approx3.0$, giving $\rho^\ast=0.0055$), so the
agreement is fortuitous at the sweep's resolution.

\subsection*{Junction myosin is keyed to the wrong variable}

The setpoint uses $\ell_e/\langle\ell\rangle$: the state variable is \emph{length}, not tension, and the
docstring's identification of the two is false in this energy, where edge tension is
$\Lambda m_e+\sum_f 2K_P(P_f-P_f^0)$ and length does not appear directly. The sign is also inverted
relative to the papers we cite: Bertet et~al.\ (2004) report myosin enriched in \emph{disassembling}
junctions, and Fernandez-Gonzalez et~al.\ (2009) find more myosin on linked edges \emph{regardless of
edge length}. The modelling literature keys on tension or strain rate and never on length
(Sknepnek et~al.\ 2023 state verbatim that their active force ``does not include the junction length'').
So our feedback is a \emph{stabiliser}---longer $\to$ more myosin $\to$ shorter, homogenising---where the
biology is a \emph{destabiliser} that generates T1s.

Worse for the identifiability argument of \S9: in the $\tau_m\to0$ limit, which is the biologically
correct one, substituting the setpoint gives
\[ E_{\text{line}}=\Lambda a\sum_e\Bigl[(1-\beta)\,\ell_e+\beta\,\ell_e^2/\langle\ell\rangle\Bigr], \]
a line tension plus \emph{a harmonic edge spring of zero rest length}. Lowering the coefficient of
variation of edge lengths is the defining property of such a spring. The measurement that $\beta$ lowers
CV at constant radius is therefore tautological: any length-restoring feedback produces it, so it cannot
discriminate a correct law from an incorrect one. The observable that would is the \textbf{T1 rate}, which
the two laws move in opposite directions and which we have never measured.

And $\langle m^{ss}\rangle=a$ holds only without the \texttt{clamp\_min(0)}. At $\beta=4$ with edge-length
CV $\approx0.32$ the setpoint clamps for any edge shorter than $0.75\langle\ell\rangle$, biasing the mean
up by an estimated 16\%---so $\beta=4$ also raises effective $\Lambda$, and ``$R$ did not move, therefore
$\beta$ is a pure redistribution'' is not established. \texttt{MYOSIN\_TRACE} already records
\texttt{myo.mean()}, so this is checkable from logs on disk without rerunning anything.

Separately, $\tau_m=20$ frames is not merely unmeasured but wrong by orders of magnitude: under the time
anchor that makes $\tau_r$ defensible ($\approx10$~min/frame), $\tau_m\approx3.5$~h against a junctional
myosin FRAP half-time of $\approx6$~s (Kasza et~al.\ 2014).

\subsection*{Two more structural objections}

\paragraph{The membrane is integrated inertially, at Re $\sim10^{-10}$.} \texttt{emit: acceleration}
gives $v\mathrel{+}=dt\,a;\;x\mathrel{+}=dt\,v$ with unit mass, where the correct equation of motion is
$\gamma\dot{\mathbf x}=\mathbf F$. The entire narrative of \S5 and \S7---the undamped spring oscillating
about a moving anchor, critical damping $c=2\sqrt k$, the tracking lag, the sheet sinking through the
epithelium, the stability reversal above $k_{\text{adh}}=10^5$---is the phenomenology of an artificial
inertia. In the overdamped limit none of it exists. Those were reported as calibration facts found by
tests rather than by inspection; they are numerical facts about a term that should not be in the model.

\paragraph{Integrin binds laminin, not collagen~IV.} Integrins $\alpha3\beta1$, $\alpha6\beta4$ and
$\alpha6\beta1$ bind laminin; the collagen~IV receptor is $\alpha1\beta1$, and \emph{Drosophila} has no
collagen-binding integrin at all. So the model couples cells to the membrane through a ligand its
reference organism's integrins do not bind, while omitting the network they do. Junctional epidermolysis
bullosa---laminin-332 \emph{or} $\alpha6\beta4$ mutations---is the genetic proof that the link runs
through laminin. Töpfer et~al.\ (2022), which we cite, also finds \emph{perlecan} contributes greatly to
elongation, contradicting our description of it as filler. And a central-force network has no bending
rigidity, which for a 100--300\,nm sheet is arguably the dominant mode.

Also: \texttt{break\_strain}$=0.35$ sits below every tissue-level failure measurement found (lens capsule
108\%, declining to $\approx$40\% with age; \emph{C.~elegans} spermatheca stretching $\approx$70\% and
recoiling hundreds of times). And the literature does not support ``thinning'': basement membranes open by
\emph{regulated focal perforation}, and where imaged during growth they are maintained or thickened.

\subsection*{Citations to fix}

``Ku \& Bilder (2023) \emph{Dev.\ Cell} 58:522'' in \texttt{membrane\_ops.py} does not exist; the real
paper is Ku, Harris \& Bilder 2023, \emph{Dev Cell} 58(3):211--223, and it is about polar cells secreting
an MMP, not basement-membrane deposition. ``Eschenbruch et~al.\ (2021) \emph{Cells} 10:1979'' could not be
verified. Okuda 2018, Hu 2018, Töpfer 2022 and Villeneuve 2024 check out.

\subsection*{The work list, in the order the audit ranks it}

\begin{enumerate}
\item \textbf{Fix \texttt{lcc} and re-report by connectivity.} One line, no new simulation: compute the
component fraction unconditionally on the final frame, and add mean degree $z=2n_{\text{bonds}}/n_{\text{alive}}$
to the reported dicts with the $z\approx4$ rigidity threshold drawn on Figure~2. This settles run~74 and
either validates or kills ``thinning, not rupture'', from trajectories already on disk.
\item \textbf{Three tissue passes that decide the ovoid.} $n=1$ and $n=2$ at fixed floor---the
dose-response the literature supports---and one run at \texttt{smooth\_phi\_deg}$=90$ to test whether the
shape survives when axisymmetry is not imposed. Then quote the estimator null and rewrite the
``amount does not matter'' controls as the consequence of the normalisation that they are.
\item \textbf{Give junction myosin the right input.} Replace $\ell_e/\langle\ell\rangle$ with the per-edge
tension already available from \op{shape\_energy\_3d}'s gradient, or with the strain rate
$d\ln\ell_e/dt$. Print $\langle m\rangle$ to test the clamp bias from existing logs. Then report
\textbf{T1 rate} and hexagonal fraction against a $\Lambda$-matched control---the observables that
discriminate a homogenising length feedback from a destabilising tension one.
\end{enumerate}

Standing large item: make the epithelium feel the sheet, for runs 70 and 74 only---the two configurations
carrying real tension, and the two where a local hernia through a fragmented patch could finally appear.

\subsection*{Where the work stands after this}

Two things survive intact. The software design---operator kinds, topological keying of per-edge state,
the negative results recorded rather than buried---is not what the audit attacks. And the combination of
a discrete crosslinked membrane, secretion-coupled membrane growth, and a 3D growing, dividing epithelium
appears to be genuinely unpublished; the reviewer could find no model with all three. That novelty is
conditional on the coupling becoming two-way and the fragmentation metric working.

What does not survive in its current form: the ovoid as a stress-driven result, the membrane failure-mode
dichotomy, and the claim that per-junction myosin expresses a tension feedback.


\end{document}
